\section{The central replacement principle}
Set
\begin{equation}
\eta=\eps/128,\quad
b=\ceil{g^{-1}\ceil{\log_2\!\frac{128}{\eps\zeta}}},\quad
L=\ceil{\frac{1024(H+b+2)}{g\eps}}. \label{eq:parameters}
\end{equation}
Logarithmic ceilings require only integer comparisons. A time-$b$ interval sends every starting law to a density in $[1-\eta,1+\eta]$ relative to $\pi$.

Consider length-$L$ blocks separated by unobserved time-$b$ intervals. A selected success is $G$. Visits in a block's initial $H$ units will be charged as possible errors.

\begin{lemma}[Error proportional to success mass]\label{lem:replace}
If each block's initial density, conditional on the previously selected block history, lies in $[1-\eta,1+\eta]$ relative to $\pi$, then the first successful block index and its entire path have TV distance at most $\eta/(1-\eta)$ from independent trials of success probability $q_G$ and successful-path law $\pi_{\rm path}(\cdot\mid G)$.
\end{lemma}
\begin{proof}
Collapse an unsuccessful block to one symbol, retaining its entire path only on success. On the success-path space, the signed difference from the stationary outcome has variation norm at most $\eta q_G$. The difference in unsuccessful mass is also at most $\eta q_G$. Thus the TV distance between these collapsed outcome kernels is at most $\eta q_G$.

At every history, success has probability at least $(1-\eta)q_G$. Couple each outcome to a fresh stationary one until success or mismatch. The mismatch probability is at most
\[
\eta q_G\sum_{k\ge0}
\Prb(\text{no actual success in the first }k\text{ blocks})
\le\eta q_G\sum_{k\ge0}[1-(1-\eta)q_G]^k
=\frac{\eta}{1-\eta}.
\]
Until mismatch the successful index and path agree. Actual hidden end states may be retained; each unobserved buffer mixes every conditional end-state law. Actual blocks are not assumed independent.
\end{proof}

Charging $\eta$ rather than $\eta q_G$ per block would lose the theorem by paying an inverse rarity. The collapsed rare outcome supplies the cancellation.

\begin{lemma}[Omitted genuine events]\label{lem:omit}
The probability of a genuine first visit in a discarded prefix or separating interval before the selected success is at most
\begin{equation}
\frac{1+\eta}{1-\eta}\frac{q(H)+q(b)}{q_G}. \label{eq:omit}
\end{equation}
When $p\le1/(8L)$ this is at most $8(H+b+2)/(gL)\le\eps/128$.
\end{lemma}
\begin{proof}
At each reached cycle the prefix-hit probability is at most $(1+\eta)q(H)$. Bound the \emph{joint} event of an unsuccessful block and a buffer hit by dropping the unsuccessful restriction. A stationary block leaves a stationary endpoint, so this probability is at most $(1+\eta)q(b)$. This bound is made before observing the block; no division by the probability of an unsuccessful block is made.

Sum over reached cycles, whose expected number is at most $1/[(1-\eta)q_G]$. If an unsuccessful block has a hit after its prefix, its first hit was necessarily in the prefix; otherwise it would satisfy $G$. Thus all omitted hits are covered. Use \eqref{eq:qg}, $q(H)+q(b)\le p(H+b+2)$, and $(1+\eta)/(1-\eta)\le2$.
\end{proof}

\section{Arbitrary starts and the complementary branch}
Obtain a positive interval $[p_L,p_U]$ with ratio at most two on its good event. Before branching, enforce positive endpoints, $p_U/p_L\le2$, and a declared polynomial output-bit budget. An invalid report produces a fixed fallback sampler immediately. The more accurate calibration below supplies this same interval. These checks make the direct-branch work bound deterministic even on a failed counting event.

\subsection{Rare branch}
If $p_U\le1/(8L)$, observe a simulated original path from $\mu$ for time $b$ and return its actual event record if it hits, including a time-zero hit. On survival, the comparison inserts an unobserved buffer of length $b$, then the separated blocks.

The unnormalized surviving law at time $b$ satisfies
\[
\mu_b^{\rm surv}\le\mu e^{b\mathcal Q}\le(1+\eta)\pi.
\]
The unconditional chance of a hit in the initial hidden buffer is therefore
\begin{equation}
(1+\eta)q(b)\le2(b+1)p\le(b+1)/(4L). \label{eq:initial}
\end{equation}
This does not divide by the probability of surviving the observed initial segment. Conditional on every surviving observed history, the subsequent buffer multiplicatively mixes the first selected block.

The algorithm does not execute that buffer or the unsuccessful blocks. Instead draw independently
\[
K\sim\operatorname{Geom}_0(q_G),\qquad
W\sim\pi_{\rm path}(\cdot\mid G).
\]
If $s_W$ is the first hit inside $W$, output
\begin{equation}
R\tau_{\rm out}=2b+K(L+b)+s_W \label{eq:clock}
\end{equation}
and the terminal state and last $H$ time units of $W$. Because $s_W>H$, the complete requested window lies inside $W$.

The ideal continuous-record error is at most
\begin{equation}
E_{\rm struct}=
\frac{\eta}{1-\eta}
+\frac{1+\eta}{1-\eta}\frac{q(H)+q(b)}{q_G}
+(1+\eta)q(b)<\eps/32. \label{eq:struct}
\end{equation}
The terms are bounded respectively by $\eps/64$, $\eps/128$, and $\eps/4096$ using \eqref{eq:parameters} and \eqref{eq:initial}. Taking the finite output record cannot increase TV.

\subsection{A sharper instance certificate}
The displayed uniform recipe is not necessary. For any integer $0\le H<L$ and buffer $b$ with a certified multiplicative mixing error $\eta<1$, suppose a calibrated lower bound $\mu_L>0$ is available for the Palm acceptance mean. The same proof, using $q_G=p(1+L)\mu_G$, gives
\begin{equation}
E_{\rm struct}\le
\frac{\eta}{1-\eta}
+\frac{1+\eta}{1-\eta}\frac{H+b+2}{(L+1)\mu_L}
+(1+\eta)p_U(b+1). \label{eq:instance}
\end{equation}
Thus a smaller chosen block can be accepted by a finite observed-mean certificate; if it fails, the uniform recipe remains available. The mean interval must include implementation bias. A relative interval for $q_G$ also bounds the geometric replacement error by the Bernoulli coupling below. This proposition supplies an actual sufficient certificate for a chosen schedule; it does not identify the uniform schedule with a necessary cost.

\subsection{Complementary branch}
If $p_U>1/(8L)$, then $p_L>1/(16L)$ on the good event. Simulate until the first hit or
\begin{equation}
T=\ceil{(gp_L)^{-1}\ceil{\log_2\!\frac{128}{\eps\zeta}}}. \label{eq:T}
\end{equation}
By \eqref{eq:killed}, the timeout probability is at most $\eps/128$. On timeout return a fixed finite fallback record. The horizon is bounded by
$16Lg^{-1}\ceil{\log_2(128/(\eps\zeta))}+1$.
Thus this branch is also polynomial: it is used only when $p^{-1}\le16L$.
